\section{Literature and Evidence Spine}

The paper rests on two layers of evidence. The first is standard portfolio and derivatives
theory: Markowitz mean--variance allocation, Black--Scholes--Merton local replication,
stochastic-calculus option accounting, and factor risk models. The second is generated
empirical evidence from local OPRA/Databento-derived option panels. The separation matters
because the paper uses different standards for different claims. A theorem needs a proof.
A literature claim needs an actual paper or book. A numerical claim needs a generated
artifact.

\subsection{Portfolio and derivatives references}

The Markowitz reference point is the maximum-Sharpe and efficient-frontier problem
\citep{markowitz1952portfolio}. The derivative reference point is the local option-pricing
logic associated with Black, Scholes, and Merton
\citep{blackscholes1973,merton1973}. This paper does not need Black--Scholes to be a
globally correct pricing model for American single-name options. It needs the weaker and
standard local statement that an option mark can be expanded in spot, volatility, and time,
with a second-order spot term. That is the same mathematical reason gamma appears in the
BSM PDE; standard stochastic-calculus treatments derive this through Ito's formula and
dynamic replication \citep{shreve2004stochasticii}.

The conditional-mean reference point is the option-return literature. Expected option
returns can reflect systematic equity exposure, volatility and variance risk premia,
skewed crash-insurance payoffs, and volatility-state dependence rather than pure pricing
noise \citep{coval_shumway2001,bakshi_kapadia2003,carr_wu2009,
bollerslev_tauchen_zhou2009,broadie_chernov_johannes2009}.  The paper uses those
mechanisms as economic forecast components, then shrinks the combined signal toward zero
before optimization.

The risk-model reference point is the factor covariance decomposition. In equity risk
models, exposures, factor covariance, and idiosyncratic covariance are separated because a
raw covariance matrix is noisy and because optimizers amplify precision-matrix error. The
option-only construction uses the same structure, but the exposure matrix is made from
Greeks rather than industry, style, or macro betas. The covariance shrinkage and
estimation discipline follows the same concern emphasized in the portfolio literature:
sample covariance matrices are noisy inputs to optimization, and small covariance errors
can be magnified when the optimizer inverts the risk matrix \citep{ledoitwolf2004}.
Backtest validation is treated with the same caution: repeated search over Sharpe ratios
raises overfitting risk, so the paper reports random feasible portfolios, bootstrap
intervals, regimes, and leave-one-underlying-out tests rather than a single selected path
\citep{bailey2017pbo,lopezdeprado2018}.

Performance is reported with complementary risk-adjusted measures.  The Sharpe ratio
follows the reward-to-variability tradition of \citet{sharpe1966}.  The Sortino ratio
replaces total volatility with target downside deviation \citep{sortino_price1994}.  The
Calmar statistic uses maximum drawdown as the loss scale, so the drawdown literature is
the relevant formal reference \citep{magdon_ismail_atiya2004}.  Omega measures the ratio
of gains above a threshold to shortfalls below it \citep{keating_shadwick2002}.  The
information ratio is computed against the SPY benchmark, following the active-return
tradition in portfolio management \citep{grinold_kahn2000}.  The empirical target for
Sortino and Omega is zero monthly return, matching the paper's no-cash-alpha convention.

\subsection{Local option data}

The empirical input is the repo-local feature store, not a live network pull. The panel
uses model-implied Greeks by design, with American-option adjustments and
European-equivalent sensitivities used as the standardized risk representation. The
filters remove nonpositive marks, stale or pathological quotes, insufficient volume,
missing Greeks, and wide relative spreads.

The quality ledger reports valid IV, delta, gamma, and vega shares. For the primary
single-name underlyings used here, valid Greek coverage is roughly 95--98\%: AAPL 95.9\%,
AMZN 97.2\%, GOOGL 97.3\%, JPM 96.0\%, META 96.6\%, MSFT 95.1\%, NVDA 97.8\%, and TSLA
97.9\%. The data provenance files and operational data-source pages are recorded in the
reproducibility note rather than cited as scholarly references.

\begin{table}[H]
\centering
\resizebox{\textwidth}{!}{\begin{tabular}{ll}
\toprule
Item & Value \\
\midrule
Raw OPRA-derived equity option rows & 160,315 \\
Raw VIX option rows after filters & 103,650 \\
Monthly snapshot dates & 129 \\
Training dates & 69 \\
Test dates & 60 \\
Primary equity underlyings & AAPL, AMZN, GOOGL, JPM, META, MSFT, NVDA, TSLA \\
VIX option treatment & VX-forward Black-76 Greeks; VIX/VVIX state only \\
VIX settlement source & exact VRO/SOQ (536 rows) \\
VIX headline status & exact VRO/SOQ settlement (excluded expiries disclosed in Section 2) \\
Post-cost layers & executable-cost scenarios (mid, half, full spread) and conservative stress-cost layer \\
Broker execution status & not broker-executed live evidence \\
Rolling option bucket assets & 54 \\
Train/test split & through 2020-12-31 / after 2020-12-31 \\
\bottomrule
\end{tabular}
}
\caption{Empirical data summary. The raw row count is after the paper's liquidity,
moneyness, and Greek-validity filters.}
\label{tab:data}
\end{table}
